GPUs have become the engine of modern high-performance computing. While NVIDIA CUDA remains the industry standard, its complexity creates a high barrier to entry. OpenAI Triton, a Python-based DSL, aims to democratize GPU programming through a block-based model. While Triton excels at dense deep learning primitives, its capability to handle irregular, latency-bound workloads remains an open question.

This paper presents a workload-driven performance analysis comparing Triton and CUDA across two domains: structured, data-parallel computations (Matrix Multiplication) and irregular, control-flow intensive tasks (Finite State Machine simulation).
Our results show that for structured workloads, Triton achieves performance parity with CUDA. For irregular workloads involving pointer-chasing and divergent control flow, CUDA kernels are 10--30$\times$ faster than the best Triton implementations. We conclude that while Triton increases productivity for dense linear algebra, low-level CUDA remains essential for high-performance irregular algorithms.
